\section{The whole associated ring and module sheaves agree}

Let $G^{\#}$ be the genuine ring sheafification of the original section
graded ring. The original structure-ring sheaf morphism
$\alpha:\mathcal O_M\to G^{\#}$ is already defined by the actual
degree-zero classes. The statements here use this action directly,
before any whole stalk dictionary. They require no coherence assumption.

\begin{definition}[The actual ring sheaf as an ambient module sheaf]
\label{hir:real-domain-actual-normal-associated-ring-ambient-module-sheaf}
\uses{hir:real-domain-actual-normal-associated-ring-sheaf}
\lean{Hironaka.HironakaSmoothCentreRealDomainActualNormalAssociatedModuleRingSheaf.actualNormalAssociatedRingSheafSection,Hironaka.HironakaSmoothCentreRealDomainActualNormalAssociatedModuleRingSheaf.actualNormalAssociatedRingSheafSectionAlgebra,Hironaka.HironakaSmoothCentreRealDomainActualNormalAssociatedModuleRingSheaf.actualNormalAssociatedRingSheafModulePresheaf,Hironaka.HironakaSmoothCentreRealDomainActualNormalAssociatedModuleRingSheaf.actualNormalAssociatedRingAmbientModuleSheaf}
\leanok
The underlying additive sheaf of $G^{\#}$ becomes a genuine ambient
module sheaf by $r\cdot s=\alpha_U(r)s$ on every original open $U$.
\end{definition}
\begin{proof}
Each $\alpha_U$ is a ring homomorphism, so multiplication by its images
defines a module structure. For $U\subseteq V$, naturality of $\alpha$
and multiplicativity of ring restriction give
$(\alpha_V(r)s)|_U=\alpha_U(r|_U)s|_U$. Thus the ring restrictions are
semilinear for the actual structure-ring restrictions. The additive
presheaf is a sheaf because forgetting commutative ring operations
preserves the limits expressing the sheaf condition. This supplies
the module sheaf with its original sections and restrictions.
\end{proof}

\begin{lemma}[The true ring sheafification unit is ambient linear]
\label{hir:real-domain-actual-normal-associated-ring-module-sheaf-unit}
\uses{hir:real-domain-actual-normal-associated-ring-ambient-module-sheaf,hir:real-domain-actual-normal-associated-module-ring-presheaf}
\lean{Hironaka.HironakaSmoothCentreRealDomainActualNormalAssociatedModuleRingSheaf.actualNormalAssociatedRingModuleSheafUnit}
\leanok
The original ring sheafification unit $G\to G^{\#}$ is a morphism
of ambient module presheaves, with the same map on every whole section.
\end{lemma}
\begin{proof}
The ambient ring sheaf morphism is defined as the original ambient
ring presheaf morphism followed by the ring sheafification unit.
The unit preserves multiplication, so it sends
$\operatorname{im}(r)s$ to $\alpha_U(r)$ times the unit image of $s$.
This is precisely ambient linearity. Its naturality is the naturality
of the original ring sheafification unit.
\end{proof}

\begin{theorem}[The whole module sheafification agrees with the actual ring sheaf]
\label{hir:real-domain-actual-normal-associated-module-ring-sheaf-comparison}
\uses{hir:real-domain-actual-normal-associated-ring-module-sheaf-unit,hir:real-domain-actual-normal-associated-module-ring-presheaf-comparison}
\lean{Hironaka.HironakaSmoothCentreRealDomainActualNormalAssociatedModuleRingSheaf.actualNormalAssociatedRingModuleSheafLift,Hironaka.HironakaSmoothCentreRealDomainActualNormalAssociatedModuleRingSheaf.actualNormalAssociatedRingModuleSheafLift_toSheaf,Hironaka.HironakaSmoothCentreRealDomainActualNormalAssociatedModuleRingSheaf.actualNormalAssociatedRingModuleSheafLiftIsIso,Hironaka.HironakaSmoothCentreRealDomainActualNormalAssociatedModuleRingSheaf.actualNormalAssociatedModuleRingSheafIso,Hironaka.HironakaSmoothCentreRealDomainActualNormalAssociatedModuleRingSheaf.actualNormalAssociatedRingModuleSheafLift_unit,Hironaka.HironakaSmoothCentreRealDomainActualNormalAssociatedModuleRingSheaf.actualNormalAssociatedModuleRingSheafIso_unit}
\leanok
The previously constructed associated module sheaf is isomorphic to
the actual ambient module sheaf underlying $G^{\#}$. The isomorphism
extends the original whole section comparison followed by the true
ring sheafification unit.
\end{theorem}
\begin{proof}
First apply module sheafification to the isomorphism from the original
module coproduct to the module presheaf underlying $G$. The preceding
linear unit extends uniquely from this latter presheaf to its module
sheafification by the sheafification adjunction. After forgetting
module structures, this extension is the canonical comparison from
the additive sheafification of the additive presheaf underlying $G$
to the additive sheaf underlying its commutative-ring sheafification.
The forgetful functor from commutative rings to abelian groups preserves
limits and filtered colimits. On the open-set site, it therefore
commutes with the two covering-colimit and matching-limit operations
constructing sheafification; their canonical comparison is an
isomorphism and has exactly the original unit factorization. Hence
the additive extension is an isomorphism. The forgetful functor from
module sheaves to additive sheaves reflects isomorphisms: an inverse
to an additive bijection that respects scalar multiplication also
respects scalar multiplication, by applying the original bijection.
Thus the linear extension is a module sheaf isomorphism. Composing
with the first sheafified isomorphism proves the assertion. The
adjunction unit identity gives the stated whole section factorization.
\end{proof}

\begin{corollary}[Every original degree is retained by the whole sheaf comparison]
\label{hir:real-domain-actual-normal-associated-module-ring-sheaf-degrees}
\uses{hir:real-domain-actual-normal-associated-module-ring-sheaf-comparison,hir:real-domain-actual-normal-associated-module-ring-presheaf-comparison}
\lean{Hironaka.HironakaSmoothCentreRealDomainActualNormalAssociatedModuleRingSheaf.actualNormalAssociatedModuleRingSheafIso_degree,Hironaka.HironakaSmoothCentreRealDomainActualNormalAssociatedModuleRingSheaf.actualNormalAssociatedModuleRingSheafIso_degree_apply}
\leanok
For every $q$, the whole sheaf comparison sends the original module
coproduct degree map, followed by its module sheafification unit,
to the literal original degree-$q$ section inclusion followed by
the ring sheafification unit, on every original open.
\end{corollary}
\begin{proof}
Precompose the whole section unit factorization with the categorical
degree-$q$ inclusion. The original presheaf comparison sends that
inclusion to the literal degree-$q$ section inclusion. Associativity
then gives the asserted identity of morphisms. Evaluating on any
original open and any original layer section gives its stated
section formula.
\end{proof}
